%! TEX root = thesis.tex

\chapter{Fundamentals}
\label{sec:fundamentals}

\section{LLM Pretraining}
\label{sec:llm-pretraining}

Large language models are neural networks with hundreds of billions of
parameters trained on trillions of text tokens.  The pretraining phase
dominates the total energy budget of an LLM lifecycle: it runs for weeks to
months on large GPU clusters.  For example, DeepSeek~V3, a 671B-parameter
mixture-of-experts model, required 2.79M GPU-hours on 2048 NVIDIA H800
GPUs~\cite{deepseek-ai_deepseek-v3_2024}.  Kimi~K2, a 1T-parameter mixture-of-experts model,
was trained on 15.5T tokens with a minimum 256-GPU model-parallel
group~\cite{kimi_team_kimi_2025}.  The effective energy consumption depends on
the GPU count, per-GPU power draw, and data center power usage effectiveness.

The training process is iterative: each step processes a batch of tokens,
computes gradients, and updates model weights.  This inherent repetitiveness
makes training amenable to pausing at natural breakpoints such as epoch
boundaries or checkpoints.  Modern training frameworks support asynchronous
checkpointing, which reduces the time overhead of saving and restoring model
state.

\section{Grid Carbon Intensity}
\label{sec:grid-carbon-intensity}

The carbon intensity of electricity measures the mass of CO$_2$ emitted per
unit of electrical energy consumed, typically expressed in
gCO$_2$eq/kWh.  Grid carbon intensity varies on multiple timescales:
\begin{itemize}
  \item \textbf{Hourly}: solar generation causes midday dips on sunny days;
    wind patterns create overnight fluctuations.
  \item \textbf{Daily}: weekday vs.\ weekend demand patterns shift the generation
    mix.
  \item \textbf{Seasonal}: winter heating demand and summer cooling demand
    alter the marginal generator.
  \item \textbf{Regional}: grids with high renewable penetration (e.g., Sweden)
    have low mean intensity; grids reliant on coal (e.g., China) have high
    mean intensity.
\end{itemize}

These variations create opportunities for temporal workload shifting.  A
training job running in Germany in 2025 experiences a mean intensity of
approximately 380~gCO$_2$eq/kWh with a standard deviation of 142~gCO$_2$eq/kWh
(coefficient of variation 0.37).  This wide range means periods of
near-zero-carbon power alternate with fossil-dominated intervals, making
Germany a promising candidate for carbon-aware scheduling.  In contrast, the
Chinese grid averages 519~gCO$_2$eq/kWh with low variability (coefficient of variation 0.11),
offering few clean intervals worth waiting for.

Carbon intensity data is available from services such as Electricity Maps,
which provides historical and fore-casted marginal intensity at 5-minute
resolution for most global regions.  Marginal intensity (the intensity of the
generator on the margin) is more appropriate for assessing the impact of an
additional load than average intensity~\cite{dodge_measuring_2022}.

\section{Temporal Workload Shifting}
\label{sec:temporal-workload-shifting}

Temporal workload shifting delays or pauses computational workloads during
periods of high grid carbon intensity and resumes them when the grid is
cleaner.  The core trade-off is between carbon savings (reduced operational
emissions) and time overhead (extended wall-clock duration).  The approach
requires no hardware modification and is applicable to any interruptible
workload with checkpointing support.

Prior work has applied temporal shifting to generic cloud batch
jobs~\cite{wiesner_lets_2021}, achieving 3--34\% operational carbon savings
depending on the delay tolerance and regional grid characteristics.  For LLM
training, the savings potential depends on model scale (larger GPU clusters
consume more power per unit time, increasing the absolute benefit of pausing),
checkpoint overhead (frequent saving and restoring adds time and energy), and
grid variability (high-variance grids offer more arbitrage opportunities).

\section{Hysteresis}
\label{sec:hysteresis}

A hysteresis policy uses separate thresholds for entering and leaving a state
to prevent rapid toggling.  In the context of carbon-aware training, two
thresholds define the policy:
\begin{itemize}
  \item $\theta_p$: the pause threshold.  When training is active and carbon
    intensity exceeds $\theta_p$, training pauses.
  \item $\theta_r$: the resume threshold.  When paused and carbon intensity
    falls below $\theta_r$, training resumes.
\end{itemize}

The condition $\theta_r \le \theta_p$ must hold.  The difference $\Delta_\theta
= \theta_p - \theta_r$ is the hysteresis margin.  A margin of zero
($\theta_r = \theta_p$) means the job toggles at a single threshold.  For a
positive margin, intensities in the interval $[\theta_r, \theta_p]$ form a dead
zone where the job remains in its current state, preventing flutter---rapid
pause-resume cycles triggered by small fluctuations around a single threshold.

Hysteresis introduces a trade-off: a wider margin reduces the number of
transitions (saving checkpoint overhead) but also leaves the job paused through
moderately clean intervals that could have been exploited productively.
Identifying the optimal margin for a given region and model is the central
focus of this work.
